\section{Background}
\label{sec:background}

\subsection{Agentic AI Architecture}

Modern AI coding assistants operate as agentic systems: the model
receives a set of tool definitions (JSON schemas describing available
operations), a system prompt (identity, instructions, behavioral
constraints), and a growing message history (alternating user and
assistant turns including tool invocations and their results).

On each API call, the client application assembles the full context:
system prompt, tool definitions, and complete message history. The
inference provider tokenizes this context, computes attention over
all tokens, and generates a response. The response is appended to
the message history, and the cycle repeats.

The critical property: \emph{context is reassembled from scratch on
every API call.} There is no persistent state at the inference layer
between calls. The entire conversation---including tool definitions
that haven't changed and results that will never be referenced
again---is sent, tokenized, and attended to on every turn.

\subsection{Context Windows as Physical Memory}

The analogy between context windows and physical memory is structural,
not merely metaphorical (Table~\ref{tab:vm-analogy}).

\begin{table}[ht]
\centering
\caption{Virtual memory analogy for context window management.}
\label{tab:vm-analogy}
\small
\begin{tabular}{@{}ll@{}}
\toprule
\textbf{OS Concept} & \textbf{Context Window Equivalent} \\
\midrule
Physical memory & Context window (200K tokens) \\
Virtual memory & Persistent state (disk, databases) \\
Page table & Retrieval handles (UUIDs, anchors) \\
Page fault & Model re-requests evicted content \\
MMU & Proxy layer between client and API \\
Working set & Currently relevant context subset \\
Demand paging & Load tool definitions on first use \\
Thrashing & Evict/fault cycle exceeding useful work \\
Pinning & Fault history prevents re-eviction \\
\bottomrule
\end{tabular}
\end{table}

Early computer systems managed limited physical memory through
manual overlays: programmers explicitly swapped code segments in and
out of a fixed address space. The transition to virtual
memory---Atlas~\citep{kilburn1962one}, Multics, then
Unix---automated this management through hardware page tables,
demand paging, and working-set-based eviction
policies~\citep{denning1968working,denning1970virtual}.

Context windows are in the overlay era. Applications manually
manage what fits, with no system-level support for eviction,
demand loading, or working set estimation.

\subsection{The Memory Hierarchy}
\label{sec:hierarchy}

The field's response to context pressure is to increase context
window size---from 4K to 32K to 128K to 1M tokens and beyond. This
is the equivalent of building machines with more physical RAM
instead of inventing virtual memory. It works but does not scale:
attention cost is quadratic in context length, and the $O(n^2)$ cost
explosion in long sessions means that even million-token windows
fill and degrade~\citep{zeyliger2026quadratic}.

The alternative is architectural: a managed hierarchy of levels,
each larger, slower, and cheaper than the one above it, connected by
eviction policies and fault mechanisms (Table~\ref{tab:hierarchy}).

\begin{table}[ht]
\centering
\caption{Memory hierarchy for LLM systems.}
\label{tab:hierarchy}
\small
\begin{tabular}{@{}llll@{}}
\toprule
\textbf{Level} & \textbf{Contents} & \textbf{Eviction} & \textbf{Fault} \\
\midrule
L1 & Generation window & --- & --- \\
L2 & Working set (pinned) & Pressure-based & Re-read \\
L3 & Session history & Age-based & Summary expansion \\
L4 & Cross-session memory & LRU / relevance & Retrieval query \\
Storage & Full corpus & None & Search + ingest \\
\bottomrule
\end{tabular}
\end{table}

\textbf{L1} is the active generation window---the tokens the model
attends to on the current API call. Small, fast, expensive per
token. This is what everyone calls ``the context window.''

\textbf{L2} is the working set: content the model is actively using
but that need not occupy L1 on every turn. Managed by fault-driven
pinning (Section~\ref{sec:pinning}). A page evicted from L1 that is
immediately re-requested is promoted to L2 and kept resident until
context pressure forces demotion.

\textbf{L3} is session history: earlier conversation turns,
completed tool interactions, superseded plans. Compressed with
declared losses into structured summaries. Faultable by expanding
the summary back to original content from the client's backing
store.

\textbf{L4} is cross-session persistent memory: authored
compressions (tensors), activity records, semantic indices. Survives
session death. Retrieved by graph traversal or similarity search.

\textbf{Storage} is the full corpus: every conversation transcript,
every tool output, every document. Archived, indexed, never resident
unless demanded.

Content migrates between levels based on access patterns. The hot
path stays in L1. Working-set content is demand-paged between L1 and
L2. Cold history compresses into L3. Cross-session knowledge
persists in L4. The context limit does not disappear---it becomes
the L1~cache size. The total addressable memory is unbounded. System
performance is determined by hit rates at each level, not by L1
capacity alone.

This paper implements and evaluates L1~management (eviction and
garbage collection) and L2~management (fault-driven pinning).
L3~(rolling conversation compaction) is implemented via cooperative
cleanup tags (Section~\ref{sec:l3}) but not yet evaluated at scale.
We describe the interface to L4~(persistent memory stores).

\subsection{Related Work}
\label{sec:related}

\paragraph{Context compression.}
LLMLingua~\citep{jiang2024longllmlingua} and related approaches compress
prompts by removing tokens predicted to be low-information. These
are lossy, model-dependent operations applied to content. Our
interventions are structural: we remove entire categories of waste
(unused tool schemas, duplicate content, dead results) without lossy
compression of the remaining content.

\paragraph{Agent context pruning.}
SWE-Pruner~\citep{wang2026swepruner} trains a 0.6B-parameter neural
skimmer for task-aware context pruning in SWE-bench agents, achieving
23--54\% token reduction. \citet{lindenbauer2025complexity} show that
simple observation masking halves agent cost while matching
LLM-summarization solve rates. ACON~\citep{kang2025acon} provides a
unified framework for history and observation compression, reducing
peak usage by 26--54\%.

These approaches frame the problem as prompt optimization or
compression. The key difference: they optimize for benchmark scores
on fixed task distributions. We measure production waste profiles
and implement eviction policies with an empirically measured fault rate. Their
framing is ``make the prompt smaller.'' Ours is ``manage the working
set.'' Working-set management composes with other interventions and
provides measurable fault rates.

\paragraph{Quadratic cost analyses.}
Recent work has quantified the $O(n^2)$ cost explosion in long
agentic sessions where cache reads dominate~\citep{zeyliger2026quadratic}.
Each API call reprocesses the full context; as sessions grow, the
cumulative cost grows quadratically with session length. Our
amplification factor measurement (Section~\ref{sec:amplification})
confirms this at production scale and our interventions directly
address it.

\paragraph{Prompt caching.}
Anthropic and OpenAI offer prompt caching that avoids recomputation
of static prefixes. In our corpus, 93.5\% of input tokens are cache
reads---caching is working. But cached tokens still occupy the
context window and require attention computation for every output
token generated. Caching reduces the input processing cost; it does
not reduce the attention cost or the memory pressure.

\paragraph{Context engineering.}
The emerging discipline of context
engineering~\citep{deepset2025context,willison2025context,mei2025survey}
recognizes that managing what enters context is as important as the
model itself. \citet{mei2025survey} survey the field comprehensively,
covering retrieval, processing, and management. Anthropic's own
engineering team describes practical
patterns~\citep{rajasekaran2025context}. Our work provides a specific
mechanism (proxy-layer interposition with eviction policies) and the
empirical production data that the field currently lacks.

\paragraph{Prompt-level memory management.}
MemGPT~\citep{packer2023memgpt} pioneered the OS virtual memory
analogy: a FIFO message queue with recursive summarization at
capacity, plus archival/recall databases accessible via tool calls.
The model can save information but cannot explicitly release it;
eviction is capacity-triggered (flush at 100\%), not cost-driven.
Contextual Memory Virtualisation~\citep{cmv2026} models session
history as a DAG with structurally lossless trimming that strips
mechanical bloat (raw tool outputs, base64, metadata) while
preserving conversational content verbatim---achieving up to 86\%
reduction for tool-heavy sessions. Focus~\citep{verma2026focus}
gives the agent autonomy over compression: the LLM consolidates
learnings into a persistent knowledge block and prunes raw history,
inspired by slime mold exploration. SideQuest~\citep{sidequest2026}
teaches the model to evict stale tool outputs from the KV cache via
a fine-tuned parallel reasoning thread, achieving 56--65\% peak
memory reduction---but eviction is irreversible and requires
inference engine modifications.

These systems address the architectural problem (MemGPT's VM
analogy), the structural problem (CMV's lossless trimming), the
autonomy problem (Focus's model-directed compression), or the
inference problem (SideQuest's KV cache eviction). None derives an
economic model where prompt-token retention cost scales quadratically
while fault cost scales linearly. None provides cooperative
recall/release primitives informed by that cost model. None reports
empirical fault rates on production sessions.

\paragraph{Positioning.}
Prior work addresses either the cost problem (quadratic analyses) or
the content problem (compression, pruning, RAG) but not the
structural problem: context windows are unmanaged physical memory.
Our contribution is the systems abstraction---working set, eviction
policy, fault rate, amplification factor---and the empirical evidence
that it works on production data.

% ====================================================================
